\chapter{Elementos Pós-textuais}

Os elementos pós-textuais completam o trabalho após a parte textual. Eles registram fontes, esclarecem termos, apresentam materiais complementares e preservam documentos necessários à compreensão da pesquisa.

\section{Referências}

As referências listam as fontes citadas no texto. Devem seguir a norma vigente adotada pela instituição e manter consistência de autoria, título, edição, local, editora, ano, link e data de acesso quando aplicável.

No modelo de TCC do UnDF-TeX, inclusive no formato de artigo, as referências mantêm a apresentação do TCC convencional: abertura em nova página, título centralizado sem indicativo numérico, entradas alinhadas à margem esquerda, espaçamento simples dentro de cada referência e separação uniforme entre elas. A camada bibliográfica é compartilhada pelos perfis. Essa padronização da entrega acadêmica não altera quais elementos são obrigatórios ou opcionais; uma futura submissão a revista seguirá suas próprias instruções editoriais.

O estudante deve conferir se toda fonte citada aparece nas referências e se toda referência listada foi efetivamente citada no texto, quando a orientação do curso exigir lista apenas de fontes citadas.

\section{Glossário}

O glossário reúne termos técnicos ou especializados acompanhados de definição. Ele é útil quando o trabalho utiliza vocabulário específico que pode dificultar a leitura por público não especializado.

Não deve ser usado para definir termos comuns ou para substituir explicações conceituais que deveriam aparecer no referencial teórico.

\section{Apêndices}

Apêndices são materiais elaborados pelo autor do trabalho. Podem incluir instrumentos de coleta, roteiros de entrevista, questionários, quadros analíticos, produtos educacionais, materiais de apoio, documentos produzidos na pesquisa ou dados complementares.

Apêndices devem ser citados no texto. Não se deve incluir material sem relação direta com o trabalho apenas para aumentar volume.

\section{Anexos}

Anexos são materiais não elaborados pelo autor, mas necessários à compreensão, comprovação ou contextualização do trabalho. Podem incluir leis, documentos institucionais, mapas, formulários externos, autorizações, protocolos ou imagens de fonte externa.

Assim como os apêndices, anexos devem ser citados no texto e ter relação clara com a pesquisa.

\section{Índice}

O índice é elemento opcional e pouco comum em TCCs. Quando utilizado, deve permitir localização de termos, nomes ou assuntos relevantes. Em trabalhos de graduação, normalmente o sumário e as listas são suficientes.

\caixaatencao{Material complementar não deve esconder fragilidades do texto. Se uma informação é essencial para entender o argumento, ela precisa ser discutida no corpo do trabalho, ainda que o documento completo apareça em apêndice ou anexo.}

